\begin{center}
\textsc{Gauss} to \textsc{Olbers}.\quad Göttingen, 2~September 1837.
\end{center}

\par\noindent\hspace*{\parindent}\mbox{.\ .\ .\ .\ .\ .}\quad \textsc{Gerling}, who has now completed his trigonometric measurements in Hesse, has now undertaken a further operation to determine the difference in longitude between Göttingen and Mannheim.
Signals are given on two mountains, Meisner and Feldberg. Those from the former are visible at my observatory, those from the latter at the Mannheim observatory, but both at once on an intermediate mountain near Marburg, where \textsc{Gerling} observes with a \textsc{Kessel} chronometer.
The signals are gunpowder signals by night and
\par\smallskip\noindent\hfill{\footnotesize 48\textsuperscript{*}}\hspace*{.12\linewidth}\par
\clearpage
\noindent\makebox[0pt][l]{\fontsize{13}{16}\selectfont \textbf{380}}\hfill{\small\textbf{CORRESPONDENCE WITH GERLING.}}\hfill\null\par\medskip
\noindent heliotropic ones by day.
The weather is unfavourable. Quite a number of signals of both kinds have already been observed here, but according to his last letter \textsc{Gerling} had still seen almost nothing from the Feldberg.\ \mbox{.\ .\ .\ .\ .\ .}

\begin{center}\rule{.24\linewidth}{.4pt}\end{center}
